\documentclass[a4paper, serif, 10pt]{moderncv}

%% moderncv
% style and color
\moderncvstyle{banking}
\moderncvcolor{black}

%% page layout/margins
\usepackage[top=2.5cm, bottom=2.5cm, left=1.5cm, right=1.5cm]{geometry}

\nopagenumbers{}

%% Babel config for Brazilian Portuguese language
% ref:
% - https://latex3.github.io/babel/guides/locale-portuguese.html
\usepackage[brazilian]{babel}

%% fontspec
\usepackage{fontspec}

\IfFontExistsTF{Linux Libertine}{
  \setmainfont[Ligatures={TeX, Common}, Numbers=OldStyle]{Linux Libertine}
}{}
\IfFontExistsTF{Linux Libertine O}{
  \setmainfont[Ligatures={TeX, Common}, Numbers=OldStyle]{Linux Libertine O}
}{}

%% CTeX
% CJK support, used to show CJK characters in the resume
%
% - fontset=none: disable builtin fontset but instead set the CJK font manually
% - heading=false: disable ctex heading
% - punct=kaiming: use kaiming punctuations styles for CJK
% - scheme=plain: use plain scheme, do not override \normalsize font size
% - space=auto: space settings for CJK characters
%
% ref:
% - http://ctan.mirrorcatalogs.com/language/chinese/ctex/ctex.pdf
\usepackage[UTF8, heading=false, punct=kaiming, scheme=plain, space=auto]{ctex}

\IfFontExistsTF{Noto Serif CJK SC}{
  \setCJKmainfont{Noto Serif CJK SC}
}{}
\IfFontExistsTF{Noto Sans CJK SC}{
  \setCJKsansfont{Noto Sans CJK SC}
}{}

%% line spacing
\usepackage{setspace}
\setstretch{1.125}

%% URL styling - use normal text instead of monospace
\urlstyle{same}

% Auto-underline all links
% Defer until \begin{document} so hyperref is already loaded
\AtBeginDocument{%
  \let\oldhref\href
  \renewcommand{\href}[2]{\oldhref{#1}{\underline{#2}}}%
}

%% Patch \httplink and \httpslink to handle full URLs
% The original moderncv macros always prepend http:// or https:// to URLs.
% This causes issues when the URL already has a protocol (e.g., https://example.com)
% resulting in malformed URLs like https://https://example.com.
% This patch checks if the URL already contains "://" and uses it directly if so.
% Note: We use \str_set:Nx (x-expansion) to expand macros like \@homepage first.
\ExplSyntaxOn
\str_new:N \g_yamlresume_url_str

\RenewDocumentCommand{\httpslink}{O{}m}{
  \str_set:Nx \g_yamlresume_url_str {#2}
  \str_if_in:NnTF \g_yamlresume_url_str {://}
    { \url{#2} }
    { \url{https://#2} }
}

\RenewDocumentCommand{\httplink}{O{}m}{
  \str_set:Nx \g_yamlresume_url_str {#2}
  \str_if_in:NnTF \g_yamlresume_url_str {://}
    { \url{#2} }
    { \url{http://#2} }
}
\ExplSyntaxOff

\name{Mariana Costa Almeida}{}
\title{Gerente de Sucesso do Cliente | Customer Success Manager}
\phone[mobile]{+55 11 98765-4321}
\email{mariana.almeida@email.com}
\homepage{https://marianaalmeida.com.br}

\address{Avenida Paulista, 1578, Bela Vista, São Paulo, São Paulo, Brasil, 01310-200}{}{}

\extrainfo{{\small \faLinkedin}\ \href{https://www.linkedin.com/in/marianacalmeida}{@marianacalmeida} {} {} {} • {} {} {} 
{\small \faMedium}\ \href{https://medium.com/@marianaalmeida}{@marianaalmeida}}

\begin{document}

\maketitle

\section{Informações Básicas}

\cvline{}{\begin{itemize}
\item Gerente de Sucesso do Cliente com mais de 8 anos de experiência em SaaS B2B, focada em retenção, expansão de receita e satisfação de clientes.
\item Especialista em construir relacionamentos de confiança com stakeholders, liderar equipes multidisciplinares e implementar estratégias de customer success baseadas em dados.
\item Experiência comprovada na redução de churn, aumento de NPS e identificação de oportunidades de upsell e cross-sell.
\item Fortes habilidades em análise de métricas, gestão de projetos, comunicação e resolução de problemas.
\item Inglês avançado e espanhol intermediário, com vivência em ambientes internacionais.
\end{itemize}}
\section{Formação}

\cventry{fev. de 2012–dez. de 2016}
        {Bacharelado, Administração de Empresas, Nota: 8.7}
        {Universidade de São Paulo}
        {\url{https://www.usp.br}}
        {}
        {\begin{itemize}
\item Formação sólida em gestão empresarial, com ênfase em relacionamento com clientes e estratégias de retenção.
\item Desenvolvimento de projetos práticos de consultoria para pequenas e médias empresas.
\item Participação ativa em empresa júnior, liderando projetos de customer success.
\end{itemize}
\textbf{Cursos}: Gestão de Clientes e CRM, Marketing de Serviços, Análise de Dados para Negócios, Gestão de Projetos, Comportamento Organizacional, Negociação e Comunicação}
\section{Experiência Profissional}

\cventry{mar. de 2021–Atual}
        {Senior Customer Success Manager}
        {NuvemTech}
        {\url{https://www.nuvemtech.com.br}}
        {}
        {\begin{itemize}
\item Gerencia uma carteira de 45 clientes enterprise, representando R\$ 12 milhões em receita recorrente anual (ARR).
\item Lidera uma equipe de 6 CSMs, responsável por onboarding, adoção, retenção e expansão.
\item Implementou um programa de QBRs estratégicos que aumentou o NPS de 42 para 68 em 12 meses.
\item Desenvolveu playbooks de sucesso do cliente que reduziram o churn em 22\% no primeiro ano.
\item Colabora com Produto, Vendas e Marketing para alinhar a voz do cliente à estratégia corporativa.
\end{itemize}
\textbf{Palavras-chave}: Retenção, Upsell, NPS, Gestão de Equipes, Onboarding}

\cventry{jan. de 2018–fev. de 2021}
        {Customer Success Manager}
        {DataFlow Sistemas}
        {\url{https://www.dataflowsistemas.com.br}}
        {}
        {\begin{itemize}
\item Gerenciou uma carteira de 120 clientes B2B, com foco em retenção e crescimento.
\item Conduziu análises de health score e implementou ações proativas para prevenir churn.
\item Aumentou a receita de upsell em 30\% por meio de identificação de oportunidades e negociação.
\item Realizou treinamentos e webinars para promover a adoção do produto.
\item Atuou como ponto de contato principal para escalonamentos e resolução de problemas.
\end{itemize}
\textbf{Palavras-chave}: CRM, Customer Success, Análise de Churn, QBR}

\cventry{jul. de 2016–dez. de 2017}
        {Analista de Sucesso do Cliente}
        {ConectaCRM}
        {\url{https://www.conectacrm.com.br}}
        {}
        {\begin{itemize}
\item Prestou suporte a clientes na implementação e uso do software de CRM.
\item Monitorou métricas de satisfação e elaborou relatórios para a liderança.
\item Participou da criação de materiais de autoatendimento e base de conhecimento.
\item Contribuiu para a redução do tempo médio de resposta em 35\%.
\end{itemize}
\textbf{Palavras-chave}: Atendimento, Suporte, Salesforce, Métricas}
\section{Idiomas}

\cvline{Português}{Nativo ou Bilíngue \hfill \textbf{Palavras-chave}: Nativo}
\cvline{Inglês}{Proficiência Profissional Fluente \hfill \textbf{Palavras-chave}: TOEFL 105, IELTS 7.5}
\cvline{Espanhol}{Proficiência Profissional Limitada \hfill \textbf{Palavras-chave}: DELE B1}
\section{Competências}

\cvline{Customer Success}{Especialista \hfill \textbf{Palavras-chave}: Onboarding, Retenção, Expansão, QBR, NPS}
\cvline{Ferramentas de CRM}{Avançado \hfill \textbf{Palavras-chave}: Salesforce, HubSpot, Gainsight, Totango}
\cvline{Análise de Dados}{Avançado \hfill \textbf{Palavras-chave}: Power BI, Tableau, Excel Avançado, SQL}
\cvline{Gestão de Projetos}{Avançado \hfill \textbf{Palavras-chave}: Agile, Scrum, Jira, Trello}
\cvline{Comunicação}{Especialista \hfill \textbf{Palavras-chave}: Apresentações, Negociação, Escuta Ativa, Feedback}
\section{Prêmios}

\cventry{dez. de 2022}
        {NuvemTech}
        {Prêmio Excelência em Customer Success}
        {}
        {}
        {Reconhecimento anual concedido aos profissionais com melhor desempenho em retenção e satisfação de clientes.}

\cventry{out. de 2020}
        {DataFlow Sistemas}
        {Top Performer Q3}
        {}
        {}
        {Destaque por superar metas de NPS e reduzir churn em 15\% no trimestre.}
\section{Certificados}

\cventry{jun. de 2021}
        {SuccessCOACHING}
        {Certified Customer Success Manager (CCSM)}
        {\url{https://www.successcoaching.com/certifications}}
        {}
        {}

\cventry{ago. de 2019}
        {SCRUMstudy}
        {Scrum Fundamentals Certified}
        {\url{https://www.scrumstudy.com/certifications}}
        {}
        {}
\section{Publicações}

\cventry{mar. de 2023}
        {Como reduzir churn com uma jornada de onboarding eficiente}
        {Blog Customer Success Brasil}
        {\url{https://www.customersuccessbrasil.com.br/blog/onboarding-eficiente}}
        {}
        {Artigo sobre melhores práticas para estruturar o onboarding de clientes SaaS, aumentando a retenção e o tempo de vida útil.}
\section{Projetos}

\cventry{jan. de 2022–dez. de 2022}
        {Implantação de um modelo escalável de CS para clientes de baixo touch.}
        {Programa de Customer Success Escalável}
        {\url{https://www.nuvemtech.com.br/cases/cs-escalavel}}
        {}
        {\begin{itemize}
\item Estruturou playbooks e automações para atendimento de mais de 1.000 clientes.
\item Reduziu o tempo de resposta em 40\% e aumentou o NPS em 12 pontos.
\item Utilizou Gainsight para monitorar saúde do cliente e disparar ações proativas.
\end{itemize}
\textbf{Palavras-chave}: Escalabilidade, Automação, Gainsight}

\cventry{mar. de 2019–nov. de 2019}
        {Redesenho do portal de autoatendimento para clientes B2B.}
        {Portal do Cliente 2.0}
        {\url{https://www.dataflowsistemas.com.br/portal}}
        {}
        {\begin{itemize}
\item Conduziu pesquisa com usuários e priorizou funcionalidades com base em dados.
\item Aumentou em 25\% a adoção do portal e reduziu tickets de suporte em 18\%.
\end{itemize}
\textbf{Palavras-chave}: UX, Autoatendimento, Feedback}
\section{Interesses}

\cvline{Voluntariado}{Mentoria, Educação, ONGs}
\cvline{Esportes}{Corrida, Ciclismo, Natação}
\cvline{Leitura}{Negócios, Biografias, Tecnologia}
\section{Voluntariado}

\cventry{fev. de 2020–dez. de 2023}
        {Mentora de Carreira}
        {Instituto Mente Ativa}
        {\url{https://www.institutomenteativa.org.br}}
        {}
        {\begin{itemize}
\item Mentoria voluntária para jovens em início de carreira na área de negócios e tecnologia.
\item Realização de workshops sobre comunicação, currículo e preparação para entrevistas.
\end{itemize}}
    

\cventry{jan. de 2023–Atual}
        {Membro do Comitê de Conteúdo}
        {Associação Brasileira de Customer Success}
        {\url{https://www.abcs.com.br}}
        {}
        {\begin{itemize}
\item Colaboração na curadoria de conteúdos e eventos para a comunidade de Customer Success.
\item Participação em painéis e discussões sobre tendências do mercado.
\end{itemize}}
    
\end{document}